\subsection{李群的态射}

定义 2. 设 G 和 H 为李群。G 到 H 的李群态射（若无混淆也简称 G 到 H 的态射）是一个从 G 到 H 的映射，它是群同态且解析。G 的自同构群记为 Aut(G)。

G 的恒等映射是态射。两个态射的复合是态射。如果 \(f: G \to H\) 和 \(f': H \to G\) 是互逆的两个态射，则 f 和 \(f'\) 是李群同构。

例。（1）设 \(G\) 为李群。对于所有 \(x \in G\)，\(\operatorname{Int}(x)\) 都是李群 \(G\) 的自同构。

\begin{enumerate}
\def\labelenumi{(\arabic{enumi})}
\setcounter{enumi}{1}
\item
  设 G 为李群。令 \(G^{\nu}\) 表示 G 的对立群，并赋予与 G 相同的流形结构。显然 \(G^{\nu}\) 是李群（称为 G 的对立李群），且映射 \(g \mapsto g^{-1}\) 是从李群 G 到李群 \(G^{\nu}\) 的同构。
\item
  设 G 为李群，E 为完备可赋范空间。G 在 E 上的解析线性表示（若无混淆也简称 G 在 E 上的线性表示）是从李群 G 到李群 \(\mathbf{GL}(\mathrm{E})\) 的态射，换言之，是一个解析映射 \(\pi\)，从 \(\mathbf{G}\) 到 \(\mathbf{GL}(\mathrm{E})\)，满足 \(\pi (gg^{\prime}) = \pi (g)\pi (g^{\prime})\) 对 \(g\)、\(g^{\prime}\) 属于 \(\mathbf{G}\) 成立。设 E 在 K 上有有限基 \((e_1,e_2,\ldots ,e_n)\)；令 \((e_1^*,e_2^*,\dots ,e_n^*)\) 为对偶基；令 \(\rho\) 为从群 G 到群 \(\mathbf{GL}(\mathrm{E})\) 的同态；则以下条件等价：
\end{enumerate}

\begin{enumerate}
\def\labelenumi{(\roman{enumi})}
\item
  \(\rho\) 是解析线性表示；
\item
  对所有 \(x \in \mathbf{E}\) 和 \(x' \in \mathbf{E}'\)，函数 \(g \mapsto \langle \rho(g)x, x' \rangle\) 在 \(\mathbf{G}\) 上是解析的；
\item
  对所有 \(i\) 和 \(j\)，函数 \(g \mapsto \langle \rho(g)e_i, e_j^*\rangle\) 在 G 上是解析的。蕴含关系 (i) \(\Rightarrow\) (ii) \(\Rightarrow\) (iii) 显然。另一方面，函数 \(u \mapsto \langle ue_i, e_j^*\rangle\) 在 \(\mathcal{L}(\mathrm{E})\) 上构成坐标系；因此它们限制到 \(\mathbf{GL}(\mathrm{E})\) 后在 \(\mathbf{GL}(E)\) 上构成坐标系，从而得到蕴含关系 (iii) \(\Rightarrow\) (i)。
\end{enumerate}

设 G 为实李群，E 为实完备可赋范空间，且 \(\rho\) 是从群 G 到群 \(\mathbf{GL}(\mathbf{E})\) 的同态。我们将在§8、定理 1 中看到，如果 \(\rho\) 连续（当 \(\mathbf{GL}(\mathbf{E})\) 带有由 \(\mathrm{L}(\mathbf{E})\) 上的范数导出的拓扑时），那么 \(\rho\) 是解析的。但请注意，这里的连续性概念不同于 Integration, 第 VIII 章，第 2 节，定义 1 (ii)（练习 1）中所考虑的概念。

\begin{enumerate}
\def\labelenumi{(\arabic{enumi})}
\setcounter{enumi}{3}
\tightlist
\item
  设 G 为实李群，E 为复完备可赋范空间。G 在 E 上的解析线性表示是 G 到 \(\mathbf{GL}(\mathrm{E})\) 的底层实李群的态射。
\end{enumerate}

命题 4. 设 G 和 H 为李群，f 为从群 G 到群 H 的同态。f 解析的充要条件是存在 G 的一个非空开子集 U，使得 \(f \mid U\) 是解析的。

该条件显然是必要的。设它成立。对于所有 \(x_0 \in G\)，\(f(x_0x) = f(x_0)f(x)\) 对所有 \(x \in U\) 成立，因此 \(f|_{x_0U}\) 是解析的。但各集合 \(x_0U\)（其中 \(x_0 \in G\)）构成 \(G\) 的一个开覆盖。

备注。若 \(f\) 在 \(e\) 处是浸入（相应地，在 \(e\) 处是浸没），显然 \(f\) 是浸入（相应地是浸没）。

